\subsection{Polynomials}

\subsubsection{Definition}

Let us take a little side trip to relate FPSs to polynomials. As should be
clear enough from the definitions, we can think of an FPS as a
\textquotedblleft polynomial with (potentially) infinitely many nonzero
coefficients\textquotedblright. This can be easily made precise. Indeed, we
can \textbf{define} polynomials as FPSs that have only finitely many nonzero coefficients:

\begin{definition}
\label{def.fps.pol}\textbf{(a)} An FPS $a\in K\left[  \left[  x\right]
\right]  $ is said to be a \emph{polynomial} if all but finitely many
$n\in\mathbb{N}$ satisfy $\left[  x^{n}\right]  a=0$ (that is, if all but
finitely many coefficients of $a$ are $0$). \medskip

\textbf{(b)} We let $K\left[  x\right]  $ be the set of all polynomials $a\in
K\left[  \left[  x\right]  \right]  $. This set $K\left[  x\right]  $ is a
subring of $K\left[  \left[  x\right]  \right]  $ (according to Theorem
\ref{thm.fps.pol.ring} below), and is called the \emph{univariate polynomial
ring} over $K$.
\end{definition}

For example, $2+3x+7x^{5}$ is a polynomial, whereas $1+x+x^{2}+x^{3}+\cdots$
is not (unless $K$ is a trivial ring).

The definition of a \textquotedblleft polynomial\textquotedblright\ that you
have seen in your abstract algebra course might be superficially different
from that in Definition \ref{def.fps.pol}; but it necessarily is equivalent.
(In fact, Definition \ref{def.fps.pol} \textbf{(a)} can be restated as
\textquotedblleft a polynomial means a $K$-linear combination of the monomials
$x^{0},x^{1},x^{2},\ldots$\textquotedblright, and it is clear that the
monomials $x^{0},x^{1},x^{2},\ldots$ in $K\left[  \left[  x\right]  \right]  $
are $K$-linearly independent; thus, the polynomial ring $K\left[  x\right]  $
as we have defined it in Definition \ref{def.fps.pol} \textbf{(b)} is a free
$K$-module with basis $\left(  x^{0},x^{1},x^{2},\ldots\right)  $. The same is
true for the polynomial ring $K\left[  x\right]  $ that you know from abstract
algebra. Moreover, the rules for adding, subtracting and multiplying
polynomials known from abstract algebra agree with the formulas for
$\mathbf{a}+\mathbf{b}$, $\mathbf{a}-\mathbf{b}$ and $\mathbf{a}%
\cdot\mathbf{b}$ that we gave in Definition \ref{def.fps.ops}.)

We owe a theorem:

\begin{theorem}
\label{thm.fps.pol.ring}The set $K\left[  x\right]  $ is a subring of
$K\left[  \left[  x\right]  \right]  $ (that is, it is closed under addition,
subtraction and multiplication, and contains the zero $\underline{0}$ and the
unity $\underline{1}$) and is a $K$-submodule of $K\left[  \left[  x\right]
\right]  $ (that is, it is closed under addition and scaling by elements of
$K$).
\end{theorem}

\begin{proof}
[Proof of Theorem \ref{thm.fps.pol.ring} (sketched).]This is a rather easy
exercise. The hardest part is to show that $K\left[  x\right]  $ is closed
under multiplication. But this, too, is easy: Let $a,b\in K\left[  x\right]
$. Then, all but finitely many $n\in\mathbb{N}$ satisfy $\left[  x^{n}\right]
a=0$ (since $a\in K\left[  x\right]  $). In other words, there exists a finite
subset $I$ of $\mathbb{N}$ such that%
\begin{equation}
\left[  x^{i}\right]  a=0\text{ for all }i\in\mathbb{N}\setminus I.
\label{pf.thm.fps.pol.ring.xia=0}%
\end{equation}
Similarly, there exists a finite subset $J$ of $\mathbb{N}$ such that
\begin{equation}
\left[  x^{j}\right]  b=0\text{ for all }j\in\mathbb{N}\setminus J.
\label{pf.thm.fps.pol.ring.xjb=0}%
\end{equation}
Consider these $I$ and $J$. Now, let $S$ be the subset $\left\{
i+j\ \mid\ i\in I\text{ and }j\in J\right\}  $ of $\mathbb{N}$. This set $S$
is again finite (since $I$ and $J$ are finite), and we can easily see (using
(\ref{pf.thm.fps.ring.xn(ab)=2})) that%
\[
\left[  x^{n}\right]  \left(  ab\right)  =0\text{ for all }n\in\mathbb{N}%
\setminus S
\]
\footnote{\textit{Proof.} Let $n\in\mathbb{N}\setminus S$. Thus,
$n\in\mathbb{N}$ and $n\notin S$. Now, (\ref{pf.thm.fps.ring.xn(ab)=2}) yields
$\left[  x^{n}\right]  \left(  ab\right)  =\sum_{i=0}^{n}\left[  x^{i}\right]
a\cdot\left[  x^{n-i}\right]  b$. We shall next show that each addend in this
sum is $0$.
\par
Indeed, let $i\in\left\{  0,1,\ldots,n\right\}  $ be arbitrary. Let $j=n-i$.
Thus, $n=i+j$. Hence, we cannot have $i\in I$ and $j\in J$ simultaneously
(because if we did, then we would have $n=\underbrace{i}_{\in I}%
+\underbrace{j}_{\in J}\in S$ (by the definition of $S$), which would
contradict $n\notin S$). Hence, we must have $i\notin I$ or $j\notin J$ (or
both). In the former case, we have $\left[  x^{i}\right]  a=0$ (by
(\ref{pf.thm.fps.pol.ring.xia=0}), since $i\notin I$ entails $i\in
\mathbb{N}\setminus I$). In the latter case, we have $\left[  x^{j}\right]
b=0$ (by (\ref{pf.thm.fps.pol.ring.xjb=0}), since $j\notin J$ entails
$j\in\mathbb{N}\setminus J$). Thus, in either case, at least one of the two
coefficients $\left[  x^{i}\right]  a$ and $\left[  x^{j}\right]  b$ is $0$,
so that their product $\left[  x^{i}\right]  a\cdot\left[  x^{j}\right]  b$ is
$0$. In other words, $\left[  x^{i}\right]  a\cdot\left[  x^{n-i}\right]  b$
is $0$ (since $j=n-i$).
\par
Forget that we fixed $i$. We thus have shown that $\left[  x^{i}\right]
a\cdot\left[  x^{n-i}\right]  b$ is $0$ for each $i\in\left\{  0,1,\ldots
,n\right\}  $. In other words, all addends of the sum $\sum_{i=0}^{n}\left[
x^{i}\right]  a\cdot\left[  x^{n-i}\right]  b$ are $0$. Hence, the whole sum
is $0$. In other words, $\left[  x^{n}\right]  \left(  ab\right)  =0$ (since
$\left[  x^{n}\right]  \left(  ab\right)  =\sum_{i=0}^{n}\left[  x^{i}\right]
a\cdot\left[  x^{n-i}\right]  b$), qed.}. Thus, all but finitely many
$n\in\mathbb{N}$ satisfy $\left[  x^{n}\right]  \left(  ab\right)  =0$ (since
$S$ is finite). This shows that $ab\in K\left[  x\right]  $. Hence, we have
shown that $K\left[  x\right]  $ is closed under multiplication. The remaining
claims of Theorem \ref{thm.fps.pol.ring} are similar but easier.
\end{proof}

\subsubsection{Reminders on rings and $K$-algebras}

As we now know, polynomials are just a special case of FPSs. However, they
have some features that FPSs don't have in general. The most important of
these features is \emph{substitution}. To wit, we can substitute an element of
$K$, or more generally an element of any $K$-algebra, into a polynomial (but
generally not into an FPS). Before we explain how, let us recall the notions
of rings and $K$-algebras (see \cite[Chapter 2 and \S 3.11]{23wa} for details
and more about them):

\begin{definition}
\label{def.alg.ring}The notion of a \emph{ring} (also known as a
\emph{noncommutative ring}) is defined in the exact same way as we defined the
notion of a commutative ring in Definition \ref{def.alg.commring}, except that
the \textquotedblleft Commutativity of multiplication\textquotedblright\ axiom
is removed.
\end{definition}

Examples of noncommutative rings\footnote{Note that the word \textquotedblleft
noncommutative ring\textquotedblright\ does not imply that the ring is not
commutative; it merely means that commutativity is not required. Thus, any
commutative ring is a noncommutative ring.} abound in linear algebra:

\begin{itemize}
\item For any $n\in\mathbb{N}$, the matrix ring $\mathbb{R}^{n\times n}$ (that
is, the ring of all $n\times n$-matrices with real entries) is a ring. This
ring is commutative if $n\leq1$, but not if $n>1$.

More generally, if $K$ is any ring (commutative or not), then the matrix ring
$K^{n\times n}$ is a ring for every $n\in\mathbb{N}$.

\item The ring $\mathbb{H}$ of
\href{https://en.wikipedia.org/wiki/Quaternion}{quaternions} is a ring that is
not commutative.

\item If $M$ is an abelian group, then the ring of all endomorphisms of $M$
(that is, the ring of all $\mathbb{Z}$-linear maps from $M$ to $M$) is a
noncommutative ring. (Its multiplication is composition of endomorphisms.)
\end{itemize}

Next, let us recall the notion of a $K$-algebra (\cite[\S 3.11]{23wa}). Recall
that $K$ is a fixed commutative ring.

\begin{definition}
\label{def.alg.Kalg}A $K$\emph{-algebra} is a set $A$ equipped with four maps%
\begin{align*}
\oplus &  :A\times A\rightarrow A,\\
\ominus &  :A\times A\rightarrow A,\\
\odot &  :A\times A\rightarrow A,\\
\rightharpoonup &  :K\times A\rightarrow A
\end{align*}
and two elements $\overrightarrow{0}\in A$ and $\overrightarrow{1}\in A$
satisfying the following properties:

\begin{enumerate}
\item The set $A$, equipped with the maps $\oplus$, $\ominus$ and $\odot$ and
the two elements $\overrightarrow{0}$ and $\overrightarrow{1}$, is a
(noncommutative) ring.

\item The set $A$, equipped with the maps $\oplus$, $\ominus$ and
$\rightharpoonup$ and the element $\overrightarrow{0}$, is a $K$-module.

\item We have%
\begin{equation}
\lambda\rightharpoonup\left(  a\odot b\right)  =\left(  \lambda\rightharpoonup
a\right)  \odot b=a\odot\left(  \lambda\rightharpoonup b\right)
\label{eq.def.alg.Kalg.scaleinv}%
\end{equation}
for all $\lambda\in K$ and $a,b\in A$.
\end{enumerate}

(Thus, in a nutshell, a $K$-algebra is a set $A$ that is simultaneously a ring
and a $K$-module, with the property that the ring $A$ and the $K$-module $A$
have the same addition, the same subtraction and the same zero, and satisfy
the additional compatibility property (\ref{eq.def.alg.Kalg.scaleinv}).)

Consequently, a $K$-algebra is automatically a ring and a $K$-module. Thus,
all the notations and shorthands that we have introduced for rings and for
$K$-modules will also be used for $K$-algebras. For example, if $A$ is a
$K$-algebra, then both maps $\odot:A\times A\rightarrow A$ and $\left.
\rightharpoonup\right.  :K\times A\rightarrow A$ will be denoted by $\cdot$
unless there is a risk of confusion. (There is rarely a risk of confusion,
since the two maps act on different inputs: $a\cdot b$ means $a\odot b$ if $a$
belongs to $A$, and means $a\rightharpoonup b$ if $a$ belongs to $K$. Often,
even when an element $a$ belongs to both $A$ and $K$, the elements $a\odot b$
and $a\rightharpoonup b$ are equal, so confusion cannot arise.)
\end{definition}

Examples of $K$-algebras include:

\begin{itemize}
\item the ring $K$ itself;

\item the ring $K\left[  \left[  x\right]  \right]  $ of FPSs (we have defined
the relevant maps in Definition \ref{def.fps.ops}, and claimed the relevant
properties in Theorem \ref{thm.fps.ring});

\item its subring $K\left[  x\right]  $ (all its maps are inherited from
$K\left[  \left[  x\right]  \right]  $);

\item the matrix ring $K^{n\times n}$ for each $n\in\mathbb{N}$;

\item any quotient ring of $K$ (that is, any ring of the form $K/I$ where $I$
is an ideal of $K$);

\item any commutative ring that contains $K$ as a subring.
\end{itemize}

Note that the axiom (\ref{eq.def.alg.Kalg.scaleinv}) in the definition of
$K$-algebra can be rewritten as%
\[
\lambda\left(  ab\right)  =\left(  \lambda a\right)  b=a\left(  \lambda
b\right)  \ \ \ \ \ \ \ \ \ \ \text{for all }\lambda\in K\text{ and }a,b\in A
\]
using our conventions (to write $ab$ for $a\odot b$ and to write $\lambda c$
for $\lambda\rightharpoonup c$). It says that scaling a product in $A$ by a
scalar in $\lambda\in K$ is equivalent to scaling either of its two factors by
$\lambda$.

\subsubsection{Evaluation aka substitution into polynomials}

We can now define what it means to substitute an element of a $K$-algebra into
a polynomial:

\begin{definition}
\label{def.pol.subs}Let $f\in K\left[  x\right]  $ be a polynomial. Let $A$ be
any $K$-algebra. Let $a\in A$ be any element. We then define an element
$f\left[  a\right]  \in A$ as follows:

Write $f$ in the form $f=\sum_{n\in\mathbb{N}}f_{n}x^{n}$ with $f_{0}%
,f_{1},f_{2},\ldots\in K$. (That is, $f_{n}=\left[  x^{n}\right]  f$ for each
$n\in\mathbb{N}$.) Then, set%
\[
f\left[  a\right]  :=\sum_{n\in\mathbb{N}}f_{n}a^{n}.
\]
(This sum is essentially finite, since $f$ is a polynomial.)

The element $f\left[  a\right]  $ is also denoted by $f\circ a$ or by
$f\left(  a\right)  $, and is called the \emph{value} of $f$ at $a$ (or the
\emph{evaluation} of $f$ at $a$, or the \emph{result of substituting }$a$ for
$x$ in $f$).
\end{definition}

The most popular notation for the value $f\left[  a\right]  $ that we have
just defined is $f\left(  a\right)  $. Unfortunately, this notation can lead
to ambiguities (for example, $f\left(  x+1\right)  $ could mean either the
value of $f$ at $x+1$, or the product of $f$ with $x+1$). By writing $f\left[
a\right]  $ or $f\circ a$ instead, I will avoid these ambiguities, although I
will not be fully consisting and will occasionally revert to writing $f\left(
a\right)  $ when confusion is unlikely.

For example, if $f=4x^{3}+2x+7$, then $f\left[  a\right]  =4a^{3}+2a+7$. For
another example, if $f=\left(  x+5\right)  ^{3}$, then $f\left[  a\right]
=\left(  a+5\right)  ^{3}$ (although this is not obvious; it follows from
Theorem \ref{thm.pol.eval.a+b} below). \medskip

If $f$ and $g$ are two polynomials in $K\left[  x\right]  $, then the value
$f\left[  g\right]  =f\circ g$ (this is the value of $f$ at $g$; it is
well-defined because $K\left[  x\right]  $ is a $K$-algebra) is also known as
the \emph{composition} of $f$ with $g$. We note that any polynomial $f\in
K\left[  x\right]  $ satisfies
\begin{align*}
f\left[  x\right]   &  =f\ \ \ \ \ \ \ \ \ \ \text{and}\\
f\left[  0\right]   &  =\left[  x^{0}\right]  f=\left(  \text{the constant
term of }f\right)  \ \ \ \ \ \ \ \ \ \ \text{and}\\
f\left[  1\right]   &  =\left[  x^{0}\right]  f+\left[  x^{1}\right]
f+\left[  x^{2}\right]  f+\cdots=\left(  \text{the sum of all coefficients of
}f\right)  .
\end{align*}
It is fairly common to write $f\left[  x\right]  $ instead of $f$ for a
polynomial, just to stress the fact that it is a polynomial in an
indeterminate called $x$. The equality $f\left[  x\right]  =f$ justifies this.
\medskip

Definition \ref{def.pol.subs} is rather versatile. For example, if
$f\in\mathbb{Z}\left[  x\right]  $ is a polynomial with integer coefficients,
then it allows evaluating $f$ at integers, at complex numbers, at residue
classes in $\mathbb{Z}/n$, at square matrices, at other polynomials and at
FPSs. More generally, a polynomial $f\in\mathbb{Z}\left[  x\right]  $ can be
evaluated at any element of any ring, since any ring is automatically a
$\mathbb{Z}$-algebra. Evaluating polynomials at square matrices is an
important idea in linear algebra (e.g.,
\href{https://en.wikipedia.org/wiki/Cayley-Hamilton_theorem}{the
Cayley--Hamilton theorem} is concerned with the characteristic polynomial of a
square matrix, evaluated at this matrix itself).

The following theorem surveys the most basic properties of values of polynomials:

\begin{theorem}
\label{thm.pol.eval.a+b}Let $A$ be a $K$-algebra. Let $a\in A$. Then: \medskip

\textbf{(a)} Any $f,g\in K\left[  x\right]  $ satisfy%
\[
\left(  f+g\right)  \left[  a\right]  =f\left[  a\right]  +g\left[  a\right]
\ \ \ \ \ \ \ \ \ \ \text{and}\ \ \ \ \ \ \ \ \ \ \left(  fg\right)  \left[
a\right]  =f\left[  a\right]  \cdot g\left[  a\right]  .
\]


\textbf{(b)} Any $\lambda\in K$ and $f\in K\left[  x\right]  $ satisfy%
\[
\left(  \lambda f\right)  \left[  a\right]  =\lambda\cdot f\left[  a\right]
.
\]


\textbf{(c)} Any $\lambda\in K$ satisfies $\underline{\lambda}\left[
a\right]  =\lambda\cdot1_{A}$, where $1_{A}$ is the unity of the ring $A$.
(This is often written as \textquotedblleft$\underline{\lambda}\left[
a\right]  =\lambda$\textquotedblright, but keep in mind that the
\textquotedblleft$\lambda$\textquotedblright\ on the right hand side of this
equality is understood to be \textquotedblleft coerced into $A$%
\textquotedblright, so it actually means \textquotedblleft the element of $A$
corresponding to $\lambda$\textquotedblright, which is $\lambda\cdot1_{A}$.)
\medskip

\textbf{(d)} We have $x\left[  a\right]  =a$. \medskip

\textbf{(e)} We have $x^{i}\left[  a\right]  =a^{i}$ for each $i\in\mathbb{N}%
$. \medskip

\textbf{(f)} Any $f,g\in K\left[  x\right]  $ satisfy $f\left[  g\left[
a\right]  \right]  =\left(  f\left[  g\right]  \right)  \left[  a\right]  $.
\end{theorem}

\begin{proof}
See \cite[Theorem 7.6.3]{19s} for parts \textbf{(a)}, \textbf{(b)},
\textbf{(c)}, \textbf{(d)} and \textbf{(e)}. Part \textbf{(f)} is
\cite[Proposition 7.6.14]{19s}.
\end{proof}
